\section{OpenRocket as comparison baseline}\label{app:or-context}

The four appendices that follow read QtRocket against OpenRocket: a sixteen-area feature
matrix (\cref{app:or-features}), a method-by-method comparison of the flight physics
(\cref{app:or-methods}), and a side-by-side reading of the two programs as software
(\cref{app:or-architecture}). This appendix frames all three. It states why OpenRocket is
the baseline and what that framing does and does not license, describes OpenRocket itself
in one page, and names the four philosophical differences between the projects that, more
than any difference in effort, explain the shape of what the comparisons find.

\subsection{Why this baseline, and what it is not}\label{app:or-context-why}

OpenRocket is the baseline for one reason: it is the de facto standard for what a mature
open-source model-rocket simulator offers. No other freely available tool combines a
seventeen-year development history, a published methods document with a validation chapter, a
six-degree-of-freedom flight model, and a component catalog broad enough to describe most
real airframes. Measuring against anything less would flatter QtRocket; measuring against
a commercial tool would put the methods out of reach behind closed source. The framing is
inherited from the companion Architecture Review (\srcf{docs/ArchitectureReviewLatex/}),
which pinned it precisely, and it bears restating in this document's own terms: QtRocket
is an independent project with similar goals. It is not a port of OpenRocket, not a
reimplementation, and at no point in this document is OpenRocket treated as a
specification to clone. Where the two projects agree (SI units throughout, a nose-tip
datum, a headless core beneath the GUI), they arrived by parallel reasoning; where they
disagree, the disagreement is usually deliberate and the owning body section says why.

\begin{designrule}[Baseline, not specification]
The comparison measures distance, not conformance. A row that reads \fabsent{} in the
feature matrix is a statement of how far QtRocket currently is from a mature tool's
surface, never a defect report against requirements QtRocket promised to meet; equally, an
area where QtRocket is ahead of the baseline is a local engineering fact, not a claim of
superiority over a finished product. OpenRocket calibrates the map. It does not draw it.
\end{designrule}

The two sides of every comparison also carry different kinds of evidence, and the
distinction holds throughout the appendix block. QtRocket claims cite file and line into
the pinned commit \code{254cd41}, the same evidentiary standard as the body of this
document. OpenRocket claims are researched from primary sources --- the
\code{openrocket/openrocket} repository at tag \code{release-24.12} and branch
\code{unstable} (examined 2026-07-07), the published technical documentation, and the
project's developer documentation --- and carry their sources in footnotes. They are
researched, not adversarially line-verified the way QtRocket claims are, and the
appendices that lean on them say so.

\subsection{OpenRocket in brief}\label{app:or-context-brief}

OpenRocket is a GPLv3 Java application\footnote{License and module overview:
\url{https://openrocket.readthedocs.io/en/latest/dev_guide/architecture.html}.} that began
as Sampo Niskanen's 2009 master's thesis at Helsinki University of Technology and has been
in continuous development since. The thesis, updated and extended, remains published as
the \emph{OpenRocket Technical Documentation}, version 13.05
(2013)\footnote{Sampo Niskanen, \emph{Development of an Open Source model rocket
simulation software}, Master's thesis, Helsinki University of Technology, 2009; and
\emph{OpenRocket technical documentation, for OpenRocket version 13.05}, 2013-05-10,
\url{https://openrocket.sourceforge.net/techdoc.pdf}. Both are indexed at
\url{https://openrocket.info/documentation.html}.} --- still the latest revision as of
this writing, and therefore the authoritative statement of the methods while being a
decade stale in places. The staleness is real and flagged where it matters: techdoc
\S 4.2 still describes a flat-Earth model with no Coriolis effect, whereas the current
source offers flat, spherical, and WGS84 geodesy, with Coriolis acceleration in every force
evaluation under the latter two.\footnote{\code{util/GeodeticComputationStrategy.java} on the \code{unstable}
branch of \url{https://github.com/openrocket/openrocket}.} \Cref{app:or-methods} cites the
techdoc for the methods and the source wherever the code has moved past it.

The current stable release is 24.12 under a year.month scheme; the final 24.12 GitHub
release is dated July 2025, and development happens on the repository's default
\code{unstable} branch.\footnote{\url{https://github.com/openrocket/openrocket/releases};
the final \code{release-24.12} is dated 2025-07-27.} The runtime requirement is Java 17,
enforced in code rather than by documentation: \code{Application.java} hard-codes
\code{SUPPORTED\_JRE\_VERSIONS = \{ 17 \}}.\footnote{\url{https://raw.githubusercontent.com/openrocket/openrocket/release-24.12/core/src/main/java/info/openrocket/core/startup/Application.java}.}
Structurally it is two Gradle subprojects that are also two Java Platform Module System
modules: \code{info.openrocket.core}, the headless simulation engine, file formats, and
design classes, reusable by other applications; and \code{info.openrocket.swing}, the GUI,
which depends on core. Services are wired by Google Guice dependency injection behind a
static locator.\footnote{Module split and entry point from the architecture page cited
above and the root \code{build.gradle}; \code{com.google.inject:guice:7.0.0} in
\code{core/build.gradle} on \code{unstable}.} The 24.12 surface --- quaternion 6-DOF
flight, extended Barrowman aerodynamics, event-driven recovery and staging, a roughly
two-dozen-type component catalog --- is one QtRocket does not attempt to match today, which
is precisely what makes it a useful yardstick.

\subsection{Four philosophical differences}\label{app:or-context-differences}

The companion review's central observation about this comparison is that the gap pattern
is explained less by effort than by philosophy. Four differences in engineering posture
recur in every appendix that follows; a reader who holds them in mind can predict most
verdicts before reading them.

\emph{Developmental order.} QtRocket builds and tests its seams before wiring in their
consumers; OpenRocket built features first and validated them against flight afterward.
QtRocket's Barrowman composition is implemented, unit-tested, and consumed by no
production code --- the test suite says so itself: ``nothing in production consumes this
in P2'' (\src{model/tests/AeroTests.cpp}{18}). OpenRocket's equivalent machinery shipped
inside a flying product and was then measured against flight telemetry and wind-tunnel
data in the techdoc's validation chapter, whose headline claim is CP and drag-coefficient
accuracy of roughly ten percent at subsonic speeds (techdoc ch.~6--7). The consequence,
developed in
\cref{sec:incomplete}, is that QtRocket's distance to the baseline is overwhelmingly a
consumption gap, not a construction gap.

\emph{Failure posture.} QtRocket is fail-closed where OpenRocket is warn-and-continue. A
self-intersecting QtRocket design is a hard, located failure: the mass computation refuses
to produce a number and throws with the first overlap diagnostic
(\src{model/parts/Part.cpp}{189}). OpenRocket's posture is the opposite and is built into
its event vocabulary: \code{SIM\_WARN} is a first-class flight event, recorded in the
flight data alongside ignition and apogee while the simulation
proceeds.\footnote{\code{FlightEvent.Type} in \code{simulation/FlightEvent.java},
\url{https://github.com/openrocket/openrocket/blob/unstable/core/src/main/java/info/openrocket/core/simulation/FlightEvent.java}.}
The trade is stated plainly in the body: fail-closed costs user forgiveness but guarantees
that every produced number came from a coherent model.

\emph{Placement intent versus geometric offsets.} OpenRocket positions a component by a
stored \code{(AxialMethod, offset)} pair --- top, middle, bottom, absolute, or after ---
whose arithmetic is purely geometric, computed from component lengths and never from mass
or CM.\footnote{\code{rocketcomponent/position/AxialMethod.java},
\url{https://github.com/openrocket/openrocket/blob/unstable/core/src/main/java/info/openrocket/core/rocketcomponent/position/AxialMethod.java}.}
QtRocket stores placement as physical intent: a \code{StationLink} names fractional
stations on both ends of a seam, an axial gap, and a seat kind selecting the radial
compatibility rule (\src{model/parts/Placement.h}{49}), and a single resolver turns intent
into geometry for simulator and visualizer alike (\cref{sec:placement}). The intent model
is arguably richer; it is also younger, and at the pinned commit only its simplest verb is
reachable from a front end.

\emph{Infrastructure-first engineering.} Where OpenRocket's effort visibly went into
product surface --- top-level \code{install4j/} and \code{snap/} installer trees sit
beside \code{core/} and \code{swing/} in the
repository\footnote{\url{https://github.com/openrocket/openrocket}, top-level directory
listing.} --- QtRocket's went into the engineering substrate: a five-platform CI workflow
building with \code{-Werror} tree-wide (\code{/WX} on MSVC) and flying full motor-ladder sweeps on every pull
request (\src{.github/workflows/cmake-multi-platform.yml}{29}; \cref{sec:testing}). The
same posture shows up inside the numerics, where the young project is occasionally ahead
of the mature one --- QtRocket's embedded-error adaptive RKF45 against OpenRocket's
heuristically stepped RK4 (\cref{app:or-methods}) --- and in what each ships without: one
has no headless scripting story, the other no installer.

\subsection{Reading guide}\label{app:or-context-guide}

The three comparisons are ordered from broad to deep, and each stands alone.
\Cref{app:or-features} is the sixteen-area feature matrix: for every area from the
component catalog to distribution, what OpenRocket offers and what QtRocket has, with each
QtRocket cell carrying a status verdict and citations. Start there for the shape of the
gap. \Cref{app:or-methods} compares the flight physics method by method --- dynamics,
aerodynamic coefficients, drag, atmosphere and wind, gravity, integration, motors,
recovery, and validation posture --- and is written for the reader who wants the actual
equations. \Cref{app:or-architecture} compares the two programs as software: module
structure, the design tree, positioning, mass properties, simulation control, extension
surfaces, persistence, motor data, and GUI architecture. All three inherit this appendix's
two conventions without restating them: distance, not conformance; and footnoted primary
sources on the OpenRocket side against pinned \code{file:line} citations on the QtRocket
side.
